\documentclass[10pt,letterpaper]{article}
\usepackage{cornell-notes}
\usepackage{mathtools}

\title{Chapter 14: Statistical Description of Data}
\author{}
\date{}

\setDocTitle{Chapter 14: Statistical Description of Data}
\setDocSubtitle{Numerical Methods}
\setDocVersion{v1.0}
\setDocStatus{Study Companion}
\setDocOwner{Numerical Methods Cornell Notes}
\setDocAuthor{}
\setDocDate{\today}
\setDocSummary{Cornell-format study and review notes for Chapter 14: Statistical Description of Data.}

\setCornellCollection{Numerical Methods Cornell Notes}
\setCornellUnitType{Chapter}
\setCornellUnitNumber{14}
\setCornellUnitTitle{Statistical Description of Data}
\setCornellCompanionLabel{Study and review companion}

\hypersetup{pdftitle={Chapter 14: Statistical Description of Data},pdfsubject={Cornell notes},pdfauthor={}}

\begin{document}
\maketitle

\section*{Chapter Overview}
This chapter organizes the mathematical and computational ideas behind statistical description of data. The notes preserve the numbered section sequence while emphasizing method selection, explicit assumptions, numerical reliability, cost, and verification.

\section*{Learning Objectives}
Explain the governing problem class; compare Moments of a Distribution: Mean, Variance, Skewness, and So Forth, Do Two Distributions Have the Same Means or Variances?, Are Two Distributions Different?, and Contingency Table Analysis of Two Distributions; design defensible stopping and verification checks; distinguish problem conditioning from algorithmic stability.

\section*{Prerequisites}
Probability, distributions, expectation, sampling, and elementary hypothesis testing.

\section*{Operational Workflow}
Define population, sampling unit, and null model; compute stable summaries; choose a statistic whose assumptions match the data; obtain uncertainty or a reference distribution; report effect size with the decision.

\subsection*{Formula and notation anchor}
\[
\bar{x}=\frac1n\sum_{i=1}^n x_i,\qquad s^2=\frac1{n-1}\sum_{i=1}^n(x_i-\bar{x})^2
\]
A two-pass or compensated calculation is safer than subtracting two large nearly equal moment estimates.

\begin{warnbox}{Numerical warning}
Confusing significance with importance, multiple testing, invalid independence assumptions, small expected counts, outliers, cancellation in variance, and interpreting correlation as causation.
\end{warnbox}

\section{14.0 Introduction}
\begin{cornell}
\noterow{What is the central idea?}{Statistical description separates sample summaries from population claims and attaches uncertainty to every inferential comparison.}
\noterow{How should it be used?}{For Introduction, begin from explicit inputs, outputs, preconditions, and representation choices.  Define population, sampling unit, and null model; compute stable summaries; choose a statistic whose assumptions match the data; obtain uncertainty or a reference distribution; report effect size with the decision.}
\noterow{Cost and reliability}{Analyze arithmetic work, storage, data movement, and convergence separately.  Relevant failure pressures include confusing significance with importance, multiple testing, invalid independence assumptions, small expected counts, outliers, cancellation in variance, and interpreting correlation as causation.  Use scaled tolerances and report both success criteria and diagnostic evidence.}
\noterow{Implementation checklist}{\begin{itemize}
\item Validate dimensions, domains, ordering, and structural assumptions.
\item Guard small divisors, invalid states, overflow, underflow, and nonfinite values.
\item Make mutation, workspace, indexing, and termination behavior explicit.
\item Test a simple exact case, a difficult valid case, a boundary case, and an expected failure.
\end{itemize}}
\end{cornell}
\begin{summarybox}{Section summary}
Statistical description separates sample summaries from population claims and attaches uncertainty to every inferential comparison.  Select this technique only when its assumptions match the data, and accept a result only after an independent residual, error estimate, invariant, or refinement check.
\end{summarybox}

\section{14.1 Moments of a Distribution: Mean, Variance, Skewness, and So Forth}
\begin{cornell}
\noterow{What is the central idea?}{Mean, variance, skewness, and kurtosis summarize location, scale, and shape, but their estimators have different sensitivity to outliers and cancellation.}
\noterow{How should it be used?}{For Moments of a Distribution: Mean, Variance, Skewness, and So Forth, begin from explicit inputs, outputs, preconditions, and representation choices.  Define population, sampling unit, and null model; compute stable summaries; choose a statistic whose assumptions match the data; obtain uncertainty or a reference distribution; report effect size with the decision.}
\noterow{Quantitative anchor}{\[
\mu_r=\operatorname{E}[(X-\operatorname{E}X)^r]
\]
State the dimensions, scaling, and validity assumptions before using this relation.  Check the resulting residual or invariant rather than trusting an iteration count alone.}
\noterow{Cost and reliability}{Analyze arithmetic work, storage, data movement, and convergence separately.  Relevant failure pressures include confusing significance with importance, multiple testing, invalid independence assumptions, small expected counts, outliers, cancellation in variance, and interpreting correlation as causation.  Use scaled tolerances and report both success criteria and diagnostic evidence.}
\noterow{Implementation checklist}{\begin{itemize}
\item Validate dimensions, domains, ordering, and structural assumptions.
\item Guard small divisors, invalid states, overflow, underflow, and nonfinite values.
\item Make mutation, workspace, indexing, and termination behavior explicit.
\item Test a simple exact case, a difficult valid case, a boundary case, and an expected failure.
\end{itemize}}
\end{cornell}
\begin{summarybox}{Section summary}
Mean, variance, skewness, and kurtosis summarize location, scale, and shape, but their estimators have different sensitivity to outliers and cancellation.  Select this technique only when its assumptions match the data, and accept a result only after an independent residual, error estimate, invariant, or refinement check.
\end{summarybox}

\section{14.2 Do Two Distributions Have the Same Means or Variances?}
\begin{cornell}
\noterow{What is the central idea?}{Tests of means or variances compare an observed statistic with its sampling distribution under a stated null model.}
\noterow{How should it be used?}{For Do Two Distributions Have the Same Means or Variances?, begin from explicit inputs, outputs, preconditions, and representation choices.  Define population, sampling unit, and null model; compute stable summaries; choose a statistic whose assumptions match the data; obtain uncertainty or a reference distribution; report effect size with the decision.}
\noterow{Quantitative lens}{Identify the governing equation, recurrence, probability law, or geometric predicate for Do Two Distributions Have the Same Means or Variances?.  Define every variable and scale; then derive a residual, error estimate, or invariant that can be checked independently.}
\noterow{Cost and reliability}{Analyze arithmetic work, storage, data movement, and convergence separately.  Relevant failure pressures include confusing significance with importance, multiple testing, invalid independence assumptions, small expected counts, outliers, cancellation in variance, and interpreting correlation as causation.  Use scaled tolerances and report both success criteria and diagnostic evidence.}
\noterow{Implementation checklist}{\begin{itemize}
\item Validate dimensions, domains, ordering, and structural assumptions.
\item Guard small divisors, invalid states, overflow, underflow, and nonfinite values.
\item Make mutation, workspace, indexing, and termination behavior explicit.
\item Test a simple exact case, a difficult valid case, a boundary case, and an expected failure.
\end{itemize}}
\end{cornell}
\begin{summarybox}{Section summary}
Tests of means or variances compare an observed statistic with its sampling distribution under a stated null model.  Select this technique only when its assumptions match the data, and accept a result only after an independent residual, error estimate, invariant, or refinement check.
\end{summarybox}

\section{14.3 Are Two Distributions Different?}
\begin{cornell}
\noterow{What is the central idea?}{Distributional tests detect broader differences than a mean comparison, with power depending on sample size and the alternative's shape.}
\noterow{How should it be used?}{For Are Two Distributions Different?, begin from explicit inputs, outputs, preconditions, and representation choices.  Define population, sampling unit, and null model; compute stable summaries; choose a statistic whose assumptions match the data; obtain uncertainty or a reference distribution; report effect size with the decision.}
\noterow{Quantitative lens}{Identify the governing equation, recurrence, probability law, or geometric predicate for Are Two Distributions Different?.  Define every variable and scale; then derive a residual, error estimate, or invariant that can be checked independently.}
\noterow{Cost and reliability}{Analyze arithmetic work, storage, data movement, and convergence separately.  Relevant failure pressures include confusing significance with importance, multiple testing, invalid independence assumptions, small expected counts, outliers, cancellation in variance, and interpreting correlation as causation.  Use scaled tolerances and report both success criteria and diagnostic evidence.}
\noterow{Implementation checklist}{\begin{itemize}
\item Validate dimensions, domains, ordering, and structural assumptions.
\item Guard small divisors, invalid states, overflow, underflow, and nonfinite values.
\item Make mutation, workspace, indexing, and termination behavior explicit.
\item Test a simple exact case, a difficult valid case, a boundary case, and an expected failure.
\end{itemize}}
\end{cornell}
\begin{summarybox}{Section summary}
Distributional tests detect broader differences than a mean comparison, with power depending on sample size and the alternative's shape.  Select this technique only when its assumptions match the data, and accept a result only after an independent residual, error estimate, invariant, or refinement check.
\end{summarybox}

\section{14.4 Contingency Table Analysis of Two Distributions}
\begin{cornell}
\noterow{What is the central idea?}{Counts in categorical cells support chi-square or likelihood-ratio comparisons when expected frequencies and independence assumptions are adequate.}
\noterow{How should it be used?}{For Contingency Table Analysis of Two Distributions, begin from explicit inputs, outputs, preconditions, and representation choices.  Define population, sampling unit, and null model; compute stable summaries; choose a statistic whose assumptions match the data; obtain uncertainty or a reference distribution; report effect size with the decision.}
\noterow{Quantitative lens}{Identify the governing equation, recurrence, probability law, or geometric predicate for Contingency Table Analysis of Two Distributions.  Define every variable and scale; then derive a residual, error estimate, or invariant that can be checked independently.}
\noterow{Cost and reliability}{Analyze arithmetic work, storage, data movement, and convergence separately.  Relevant failure pressures include confusing significance with importance, multiple testing, invalid independence assumptions, small expected counts, outliers, cancellation in variance, and interpreting correlation as causation.  Use scaled tolerances and report both success criteria and diagnostic evidence.}
\noterow{Implementation checklist}{\begin{itemize}
\item Validate dimensions, domains, ordering, and structural assumptions.
\item Guard small divisors, invalid states, overflow, underflow, and nonfinite values.
\item Make mutation, workspace, indexing, and termination behavior explicit.
\item Test a simple exact case, a difficult valid case, a boundary case, and an expected failure.
\end{itemize}}
\end{cornell}
\begin{summarybox}{Section summary}
Counts in categorical cells support chi-square or likelihood-ratio comparisons when expected frequencies and independence assumptions are adequate.  Select this technique only when its assumptions match the data, and accept a result only after an independent residual, error estimate, invariant, or refinement check.
\end{summarybox}

\section{14.5 Linear Correlation}
\begin{cornell}
\noterow{What is the central idea?}{Correlation normalizes covariance to measure linear association, but it is not causation and can be distorted by outliers or nonlinear structure.}
\noterow{How should it be used?}{For Linear Correlation, begin from explicit inputs, outputs, preconditions, and representation choices.  Define population, sampling unit, and null model; compute stable summaries; choose a statistic whose assumptions match the data; obtain uncertainty or a reference distribution; report effect size with the decision.}
\noterow{Quantitative anchor}{\[
R_{xy,k}=X_k^*Y_k
\]
State the dimensions, scaling, and validity assumptions before using this relation.  Check the resulting residual or invariant rather than trusting an iteration count alone.}
\noterow{Cost and reliability}{Analyze arithmetic work, storage, data movement, and convergence separately.  Relevant failure pressures include confusing significance with importance, multiple testing, invalid independence assumptions, small expected counts, outliers, cancellation in variance, and interpreting correlation as causation.  Use scaled tolerances and report both success criteria and diagnostic evidence.}
\noterow{Implementation checklist}{\begin{itemize}
\item Validate dimensions, domains, ordering, and structural assumptions.
\item Guard small divisors, invalid states, overflow, underflow, and nonfinite values.
\item Make mutation, workspace, indexing, and termination behavior explicit.
\item Test a simple exact case, a difficult valid case, a boundary case, and an expected failure.
\end{itemize}}
\end{cornell}
\begin{summarybox}{Section summary}
Correlation normalizes covariance to measure linear association, but it is not causation and can be distorted by outliers or nonlinear structure.  Select this technique only when its assumptions match the data, and accept a result only after an independent residual, error estimate, invariant, or refinement check.
\end{summarybox}

\section{14.6 Nonparametric or Rank Correlation}
\begin{cornell}
\noterow{What is the central idea?}{Rank-based association reduces dependence on marginal scales and can detect monotone relationships more robustly than raw-value correlation.}
\noterow{How should it be used?}{For Nonparametric or Rank Correlation, begin from explicit inputs, outputs, preconditions, and representation choices.  Define population, sampling unit, and null model; compute stable summaries; choose a statistic whose assumptions match the data; obtain uncertainty or a reference distribution; report effect size with the decision.}
\noterow{Quantitative anchor}{\[
R_{xy,k}=X_k^*Y_k
\]
State the dimensions, scaling, and validity assumptions before using this relation.  Check the resulting residual or invariant rather than trusting an iteration count alone.}
\noterow{Cost and reliability}{Analyze arithmetic work, storage, data movement, and convergence separately.  Relevant failure pressures include confusing significance with importance, multiple testing, invalid independence assumptions, small expected counts, outliers, cancellation in variance, and interpreting correlation as causation.  Use scaled tolerances and report both success criteria and diagnostic evidence.}
\noterow{Implementation checklist}{\begin{itemize}
\item Validate dimensions, domains, ordering, and structural assumptions.
\item Guard small divisors, invalid states, overflow, underflow, and nonfinite values.
\item Make mutation, workspace, indexing, and termination behavior explicit.
\item Test a simple exact case, a difficult valid case, a boundary case, and an expected failure.
\end{itemize}}
\end{cornell}
\begin{summarybox}{Section summary}
Rank-based association reduces dependence on marginal scales and can detect monotone relationships more robustly than raw-value correlation.  Select this technique only when its assumptions match the data, and accept a result only after an independent residual, error estimate, invariant, or refinement check.
\end{summarybox}

\section{14.7 Information-Theoretic Properties of Distributions}
\begin{cornell}
\noterow{What is the central idea?}{Entropy and divergence quantify uncertainty and distributional separation through log-probability functionals.}
\noterow{How should it be used?}{For Information-Theoretic Properties of Distributions, begin from explicit inputs, outputs, preconditions, and representation choices.  Define population, sampling unit, and null model; compute stable summaries; choose a statistic whose assumptions match the data; obtain uncertainty or a reference distribution; report effect size with the decision.}
\noterow{Quantitative lens}{Identify the governing equation, recurrence, probability law, or geometric predicate for Information-Theoretic Properties of Distributions.  Define every variable and scale; then derive a residual, error estimate, or invariant that can be checked independently.}
\noterow{Cost and reliability}{Analyze arithmetic work, storage, data movement, and convergence separately.  Relevant failure pressures include confusing significance with importance, multiple testing, invalid independence assumptions, small expected counts, outliers, cancellation in variance, and interpreting correlation as causation.  Use scaled tolerances and report both success criteria and diagnostic evidence.}
\noterow{Implementation checklist}{\begin{itemize}
\item Validate dimensions, domains, ordering, and structural assumptions.
\item Guard small divisors, invalid states, overflow, underflow, and nonfinite values.
\item Make mutation, workspace, indexing, and termination behavior explicit.
\item Test a simple exact case, a difficult valid case, a boundary case, and an expected failure.
\end{itemize}}
\end{cornell}
\begin{summarybox}{Section summary}
Entropy and divergence quantify uncertainty and distributional separation through log-probability functionals.  Select this technique only when its assumptions match the data, and accept a result only after an independent residual, error estimate, invariant, or refinement check.
\end{summarybox}

\section{14.8 Do Two-Dimensional Distributions Differ?}
\begin{cornell}
\noterow{What is the central idea?}{Multivariate two-sample comparisons require spatial statistics whose null behavior is more complex than a one-dimensional ordered test.}
\noterow{How should it be used?}{For Do Two-Dimensional Distributions Differ?, begin from explicit inputs, outputs, preconditions, and representation choices.  Define population, sampling unit, and null model; compute stable summaries; choose a statistic whose assumptions match the data; obtain uncertainty or a reference distribution; report effect size with the decision.}
\noterow{Quantitative lens}{Identify the governing equation, recurrence, probability law, or geometric predicate for Do Two-Dimensional Distributions Differ?.  Define every variable and scale; then derive a residual, error estimate, or invariant that can be checked independently.}
\noterow{Cost and reliability}{Analyze arithmetic work, storage, data movement, and convergence separately.  Relevant failure pressures include confusing significance with importance, multiple testing, invalid independence assumptions, small expected counts, outliers, cancellation in variance, and interpreting correlation as causation.  Use scaled tolerances and report both success criteria and diagnostic evidence.}
\noterow{Implementation checklist}{\begin{itemize}
\item Validate dimensions, domains, ordering, and structural assumptions.
\item Guard small divisors, invalid states, overflow, underflow, and nonfinite values.
\item Make mutation, workspace, indexing, and termination behavior explicit.
\item Test a simple exact case, a difficult valid case, a boundary case, and an expected failure.
\end{itemize}}
\end{cornell}
\begin{summarybox}{Section summary}
Multivariate two-sample comparisons require spatial statistics whose null behavior is more complex than a one-dimensional ordered test.  Select this technique only when its assumptions match the data, and accept a result only after an independent residual, error estimate, invariant, or refinement check.
\end{summarybox}

\section{14.9 Savitzky-Golay Smoothing Filters}
\begin{cornell}
\noterow{What is the central idea?}{Local polynomial least squares produce convolution coefficients that smooth noise while preserving low-order trends and derivative information.}
\noterow{How should it be used?}{For Savitzky-Golay Smoothing Filters, begin from explicit inputs, outputs, preconditions, and representation choices.  Define population, sampling unit, and null model; compute stable summaries; choose a statistic whose assumptions match the data; obtain uncertainty or a reference distribution; report effect size with the decision.}
\noterow{Quantitative lens}{Identify the governing equation, recurrence, probability law, or geometric predicate for Savitzky-Golay Smoothing Filters.  Define every variable and scale; then derive a residual, error estimate, or invariant that can be checked independently.}
\noterow{Cost and reliability}{Analyze arithmetic work, storage, data movement, and convergence separately.  Relevant failure pressures include confusing significance with importance, multiple testing, invalid independence assumptions, small expected counts, outliers, cancellation in variance, and interpreting correlation as causation.  Use scaled tolerances and report both success criteria and diagnostic evidence.}
\noterow{Implementation checklist}{\begin{itemize}
\item Validate dimensions, domains, ordering, and structural assumptions.
\item Guard small divisors, invalid states, overflow, underflow, and nonfinite values.
\item Make mutation, workspace, indexing, and termination behavior explicit.
\item Test a simple exact case, a difficult valid case, a boundary case, and an expected failure.
\end{itemize}}
\end{cornell}
\begin{summarybox}{Section summary}
Local polynomial least squares produce convolution coefficients that smooth noise while preserving low-order trends and derivative information.  Select this technique only when its assumptions match the data, and accept a result only after an independent residual, error estimate, invariant, or refinement check.
\end{summarybox}

\section*{Method comparison}
\begin{center}
\small
\begin{tabularx}{\linewidth}{>{\bfseries\raggedright\arraybackslash}p{0.22\linewidth}XX}
\toprule
Method or lens & Best use & Main caution \\
\midrule
Parametric test & A justified distribution model & Efficient when assumptions hold \\
Rank test & Ordinal or nonnormal data & Tests a different population feature \\
Resampling & Complex statistic with representative data & Computation-heavy; sampling design remains decisive \\
\bottomrule
\end{tabularx}
\end{center}

\section*{Worked example}
\begin{exambox}{Independent worked example}
For $x=[2,4,4,6]$, $\bar{x}=4$ and $s^2=8/3$.  Centered deviations sum to zero, providing a basic arithmetic check.  The sample is too small to justify a distributional conclusion; report uncertainty and inspect the data rather than relying on one statistic.
\end{exambox}

\section*{Chapter synthesis}
\begin{summarybox}{Chapter synthesis}
The chapter's methods should be viewed as a decision system rather than a list of routines.  Begin with the mathematical problem and its conditioning, exploit structure, select an algorithm with compatible stability and cost, and verify the computed answer with evidence that is independent of the internal iteration.  The recurring implementation discipline is: define population, sampling unit, and null model; compute stable summaries; choose a statistic whose assumptions match the data; obtain uncertainty or a reference distribution; report effect size with the decision.
\end{summarybox}

\subsection*{Common mistakes and numerical hazards}
\begin{warnbox}{Common mistakes and numerical hazards}
Confusing significance with importance, multiple testing, invalid independence assumptions, small expected counts, outliers, cancellation in variance, and interpreting correlation as causation.  A second recurring mistake is reporting only a computed value: always retain tolerances, iteration counts, residuals, refinement behavior, and relevant structural checks.
\end{warnbox}

\subsection*{Retrieval practice}
\textbf{Questions}
\begin{enumerate}
\item What problem does Moments of a Distribution: Mean, Variance, Skewness, and So Forth address, and which assumption is easiest to violate?
\item What problem does Do Two Distributions Have the Same Means or Variances? address, and which assumption is easiest to violate?
\item What problem does Are Two Distributions Different? address, and which assumption is easiest to violate?
\item What problem does Do Two-Dimensional Distributions Differ? address, and which assumption is easiest to violate?
\item What problem does Savitzky-Golay Smoothing Filters address, and which assumption is easiest to violate?
\item How do conditioning, stability, and the final residual play different roles in this chapter?
\end{enumerate}

\textbf{Brief answers}
\begin{enumerate}
\item Mean, variance, skewness, and kurtosis summarize location, scale, and shape, but their estimators have different sensitivity to outliers and cancellation. Recheck the section's domain, structure, scaling, and failure diagnostics.
\item Tests of means or variances compare an observed statistic with its sampling distribution under a stated null model. Recheck the section's domain, structure, scaling, and failure diagnostics.
\item Distributional tests detect broader differences than a mean comparison, with power depending on sample size and the alternative's shape. Recheck the section's domain, structure, scaling, and failure diagnostics.
\item Multivariate two-sample comparisons require spatial statistics whose null behavior is more complex than a one-dimensional ordered test. Recheck the section's domain, structure, scaling, and failure diagnostics.
\item Local polynomial least squares produce convolution coefficients that smooth noise while preserving low-order trends and derivative information. Recheck the section's domain, structure, scaling, and failure diagnostics.
\item Conditioning describes sensitivity of the mathematical problem, stability describes error propagation by the algorithm, and the residual measures satisfaction of the computed equations; none is a universal substitute for the others.
\end{enumerate}

\subsection*{Mastery checklist}
\begin{itemize}
\item[$\square$] I can explain and select Moments of a Distribution: Mean, Variance, Skewness, and So Forth without relying on a memorized code listing.
\item[$\square$] I can explain and select Do Two Distributions Have the Same Means or Variances? without relying on a memorized code listing.
\item[$\square$] I can explain and select Are Two Distributions Different? without relying on a memorized code listing.
\item[$\square$] I can explain and select Contingency Table Analysis of Two Distributions without relying on a memorized code listing.
\item[$\square$] I can explain and select Linear Correlation without relying on a memorized code listing.
\item[$\square$] I can state a scaled stopping rule and an independent verification check.
\item[$\square$] I can identify at least one conditioning risk and one algorithmic stability risk.
\item[$\square$] I can estimate time, storage, and data-movement costs at the appropriate level.
\end{itemize}

\end{document}

